\section*{अध्याय \ref{ch:g11:fundamental-interactions} --- मूलभूत अंतःक्रियाएँ}

\begin{solution}{exo:g11:fundamental-interactions:1}
प्रोटॉन: $10^{-15}$ (\omterm{def:g11:fundamental-interactions:ladder}{न्यूक्लियॉन} की मंज़िल); पानी का \omterm{def:g11:fundamental-interactions:ladder}{अणु}: $10^{-10}$
(\omterm{def:g11:fundamental-interactions:ladder}{अणु} की मंज़िल --- अकेले \omterm{def:g11:fundamental-interactions:ladder}{अणु} \omterm{def:g11:fundamental-interactions:ladder}{परमाणु} जितने आकार के होते हैं);
रेत: $10^{-4}$ (रोज़मर्रा का पदार्थ); पृथ्वी: $10^{7}$ (ग्रह); आकाशगंगा: $10^{21}$
(गैलेक्सी)। सब \omterm{def:g10:orders-of-magnitude:unit}{मीटर} में।
\end{solution}

\begin{solution}{exo:g11:fundamental-interactions:2}
ऋणात्मक आवेश: उसने $N = \num{4.8e-9}/\num{1.6e-19} = \num{3.0e10}$ इलेक्ट्रॉन \emph{पाए} हैं। आवेश के संरक्षण से
(\cref{prop:g11:fundamental-interactions:conservation}) बालों पर $+\qty{4.8e-9}{C}$ रहता है।
\end{solution}

\begin{solution}{exo:g11:fundamental-interactions:3}
$F = \num{8.99e9} \times \dfrac{\num{2.0e-6} \times \num{3.0e-6}}{0.30^2}
= \dfrac{\num{5.39e-2}}{0.090} \approx \qty{0.60}{N}$, आकर्षी (चिह्न विपरीत हैं)। दोनों बलों का परिमाण एक, दिशाएँ
विपरीत।
\end{solution}

\begin{solution}{exo:g11:fundamental-interactions:4}
$F_0/4$; \ $9F_0$; \ $2F_0$; \ $F_0$ (अंश $\times 4$, हर $\times 4$)।
\end{solution}

\begin{solution}{exo:g11:fundamental-interactions:5}
(क) \omterm{def:g11:fundamental-interactions:four}{गुरुत्वाकर्षण}; (ख) विद्युत्-चुंबकीय; (ग) प्रबल; (घ) दुर्बल;
(ङ) विद्युत्-चुंबकीय।
\end{solution}

\begin{solution}{exo:g11:fundamental-interactions:6}
$F_E = \num{8.99e9} \times (\num{1.60e-19})^2/(\num{2.0e-15})^2
= \num{2.30e-28}/\num{4.0e-30} \approx \qty{58}{N}$; $F_G = \num{6.67e-11} \times (\num{1.67e-27})^2/\num{4.0e-30}
\approx \qty{4.7e-35}{N}$। अनुपात $F_E/F_G \approx \num{1.2e36}$। नाभिक को \omterm{def:g11:fundamental-interactions:strong}{प्रबल अंतःक्रिया}
थामे रखती है, गुरुत्व का यहाँ कोई मतलब नहीं।
\end{solution}

\begin{solution}{exo:g11:fundamental-interactions:7}
$F_E = \num{2.30e-28}/(\num{5.3e-11})^2 = \num{2.30e-28}/\num{2.81e-21}
\approx \qty{8.2e-8}{N}$; $F_G = \num{6.67e-11} \times \num{1.67e-27} \times \num{9.11e-31} /
\num{2.81e-21} \approx \qty{3.6e-47}{N}$। अनुपात $\approx \num{2.3e39}$: परमाणुओं को
\omterm{def:g11:fundamental-interactions:four}{विद्युत्-चुंबकीय अंतःक्रिया} जोड़े रखती है।
\end{solution}

\begin{solution}{exo:g11:fundamental-interactions:8}
संपर्क से आवेश बँट जाता है: $q_A = q_B = \qty{+4.0e-9}{C}$ (कुल अब भी $\qty{8.0e-9}{C}$)। फिर $F = \num{8.99e9} \times (\num{4.0e-9})^2/0.10^2 =
\qty{1.4e-5}{N}$,
प्रतिकर्षी (चिह्न समान हैं)।
\end{solution}

\begin{solution}{exo:g11:fundamental-interactions:9}
$d = \sqrt{k q^2/F} = \sqrt{\num{8.99e9} \times \num{1.0e-12}/0.90}
\approx \qty{0.10}{m}$। $9$ गुना छोटे बल के लिए $3$ गुना बड़ी दूरी चाहिए:
$d = \qty{0.30}{m}$।
\end{solution}

\begin{solution}{exo:g11:fundamental-interactions:10}
$F = \num{2.30e-28}/(\num{1.0e-10})^2 = \qty{2.3e-8}{N}$। \omterm{def:g11:fundamental-interactions:strong}{प्रबल अंतःक्रिया} का कोई योगदान नहीं: $10^{-10}$~m उसके
\omterm{def:g11:fundamental-interactions:strong}{परास} का $10^{5}$ गुना है। अणुओं को \omterm{def:g11:fundamental-interactions:four}{विद्युत्-चुंबकीय
अंतःक्रिया} बाँधती है।
\end{solution}

\begin{solution}{exo:g11:fundamental-interactions:11}
हर सिक्के पर $q = \num{e23}/\num{e6} \times \num{1.6e-19} = \qty{1.6e-2}{C}$। $F = \num{8.99e9} \times (\num{1.6e-2})^2/1.0^2 \approx
\qty{2.3e6}{N}$ --- यानी $\num{2.3e6}/9.81 \approx
\qty{2.3e5}{kg}$ का भार, एक लदी हुई
मालगाड़ी, और वह भी दस लाखवें हिस्से भर के इलेक्ट्रॉनों से। रोज़मर्रा का
पदार्थ दस लाख में एक भाग से कहीं बेहतर उदासीन होना ही चाहिए --- और है।
\end{solution}

\begin{solution}{exo:g11:fundamental-interactions:12}
$F = \num{2.30e-28}/(\num{1.4e-14})^2 = \num{2.30e-28}/\num{1.96e-28}
\approx \qty{1.2}{N}$ --- इस पैमाने पर भी बहुत बड़ा। नहीं: वे प्रबल \omterm{def:g11:fundamental-interactions:strong}{परास} से $14$
गुना दूर हैं; हर प्रोटॉन केवल अपने निकटतम पड़ोसियों से चिपका है। युग्म:
$\frac{92 \times 91}{2} = 4186$, और \emph{सभी} किसी भी दूरी पर प्रतिकर्षित करते हैं, जबकि गोंद आकार
के साथ नहीं बढ़ती। न्यूट्रॉन प्रतिकर्षण जोड़े बिना आकर्षण जोड़ते हैं और
प्रोटॉनों को दूर-दूर कर देते हैं; पर प्रतिकर्षण प्रोटॉन-संख्या के वर्ग के
साथ बढ़ता है और गोंद केवल पड़ोसियों की संख्या के साथ, इसलिए लगभग $80$
प्रोटॉन से आगे कोई भी न्यूट्रॉन-बजट हिसाब नहीं बिठा पाता --- ऐसे नाभिक क्षय
कर जाते हैं (\cref{ch:g11:nucleus-radioactivity})।
\end{solution}

\begin{solution}{exo:g11:fundamental-interactions:13}
साम्यावस्था: \omterm{def:g10:universal-gravitation:weight}{ऊर्ध्वाधर} दिशा में $T\cos\theta = mg$, क्षैतिज दिशा में $T\sin\theta = F$,
इसलिए $F = mg\tan\theta$। यहाँ $\sin\theta = 3.0/50 = 0.060$, और छोटे कोणों के लिए $\tan\theta \approx \sin\theta = 0.060$। फिर $F = \num{1.0e-3} \times 9.81 \times 0.060 \approx \qty{5.9e-4}{N}$; $q = d\sqrt{F/k} = 0.060 \times \sqrt{\num{5.9e-4}/\num{8.99e9}}
\approx \qty{1.54e-8}{C}$;
$N = \num{1.54e-8}/\num{1.6e-19} \approx
\num{9.6e10}$ \omterm{def:g11:fundamental-interactions:elementary}{मूल आवेश}।
\end{solution}

\begin{solution}{exo:g11:fundamental-interactions:14}
$F = \num{6.67e-11} \times \num{5.97e24} \times \num{7.35e22} /
(\num{3.84e8})^2 = \num{2.93e37}/\num{1.47e17} \approx \qty{2.0e20}{N}$। $kQ^2/d^2 = F$ रखने पर: $Q = d\sqrt{F/k} = \num{3.84e8} \times \sqrt{\num{2.0e20}/\num{8.99e9}}
\approx \qty{5.7e13}{C}$। यह $N = \num{5.7e13}/\num{1.6e-19} \approx \num{3.6e32}$ इलेक्ट्रॉन हैं, जिनका द्रव्यमान
$\num{3.6e32} \times \num{9.11e-31} \approx \qty{3.2e2}{kg}$ है: चंद्रमा के $\qty{7.35e22}{kg}$ का काम कुछ सौ \omterm{def:g10:orders-of-magnitude:unit}{किलोग्राम} इलेक्ट्रॉन कर
सकते हैं --- एक ही संख्या में बिजली की ताक़त और गुरुत्व की कमज़ोरी।
\end{solution}

\begin{solution}{exo:g11:fundamental-interactions:15}
$F = \num{2.30e-28}/(\num{2.8e-10})^2 = \num{2.30e-28}/\num{7.84e-20}
\approx \qty{2.9e-9}{N}$; भार $P = \num{5.9e-26} \times 9.81 \approx \qty{5.8e-25}{N}$; अनुपात $F/P \approx \num{5e15}$। हर बंध की ताक़त \emph{नियत} है, पर
भार ऊपर के पूरे द्रव्यमान के साथ बढ़ता जाता है: चट्टान पर चट्टान
चढ़ाइए और सबसे नीचे की परत पर बोझ बेहिसाब बढ़ता है जबकि उसके बंध वही रहते
हैं। लगभग \qty{e4}{m} चट्टान पर बोझ विद्युत बंधों को कुचल देता है --- गुरुत्व
प्रति कण की ताक़त से नहीं, बल्कि आवेश के कट जाने पर \emph{जुड़ते जाने} से
जीतता है।
\end{solution}

\begin{solution}{pb:g11:fundamental-interactions:1}
\textbf{1.} \omterm{def:g11:fundamental-interactions:ladder}{क्वार्क} $< \qty{e-18}{m}$; \omterm{def:g11:fundamental-interactions:ladder}{न्यूक्लियॉन} $\qty{e-15}{m}$;
नाभिक $10^{-15}$--$\qty{e-14}{m}$; \omterm{def:g11:fundamental-interactions:ladder}{परमाणु} $\qty{e-10}{m}$; \omterm{def:g11:fundamental-interactions:ladder}{अणु}
$\qty{e-9}{m}$; रोज़मर्रा की वस्तु $\qty{1}{m}$; ग्रह $\qty{e7}{m}$; ग्रह-मंडल $\qty{e11}{m}$;
आकाशगंगा $\qty{e21}{m}$।

\textbf{2.} पैमाना-गुणक $10^{5}$: \omterm{def:g11:fundamental-interactions:ladder}{परमाणु} $\qty{1}{cm} \times 10^{5} = \qty{1}{km}$ चौड़ा होगा।

\textbf{3.} $(10^{-15}/10^{-10})^3 = 10^{-15}$: पदार्थ $\num{99.9999999999999}\,\%$ ख़ाली है।

\textbf{4.} $\num{7e-5}/\num{e-10} = \num{7e5}$ \omterm{def:g11:fundamental-interactions:ladder}{परमाणु} --- एक बाल की चौड़ाई में लगभग दस लाख।

\textbf{5.} $10^{-18}$ से $10^{21}$ तक: दस की $39$ घातें।

\textbf{6.} $F_E = \num{8.99e9} \times (\num{1.60e-19})^2 /
(\num{5.3e-11})^2 \approx \qty{8.2e-8}{N}$।

\textbf{7.} $F_G = \num{6.67e-11} \times \num{1.67e-27} \times
\num{9.11e-31}/(\num{5.3e-11})^2 \approx \qty{3.6e-47}{N}$; अनुपात $F_E/F_G \approx \num{2.3e39}$। \omterm{def:g11:fundamental-interactions:ladder}{परमाणु} को बिजली थामे रखती है;
गुरुत्व केवल तमाशबीन है।

\textbf{8.} गुरुत्व के बिना: नापने लायक़ कुछ नहीं बदलता (वह $10^{39}$ नीचे है)।
बिजली के बिना: न \omterm{def:g11:fundamental-interactions:ladder}{परमाणु}, न \omterm{def:g11:fundamental-interactions:ladder}{अणु}, पदार्थ ही नहीं।

\textbf{9.} ``संस्पर्श'' असल में $\qty{e-10}{m}$ तक पास लाए गए परमाणुओं के
इलेक्ट्रॉन-बादलों का विद्युत प्रतिकर्षण है: कुछ भी कभी नहीं छूता। फ़र्श की
प्रतिक्रिया और \omterm{def:g10:inertia:inventory}{घर्षण} थोक में \omterm{def:g11:fundamental-interactions:charge}{कूलॉम} का नियम ही हैं।

\textbf{10.} हर एक पर $q = \num{e23} \times \num{e-9} \times \num{1.6e-19} =
\qty{1.6e-5}{C}$; $F = \num{8.99e9} \times (\num{1.6e-5})^2 \approx \qty{2.3}{N}$ --- एक अरब में से एक इलेक्ट्रॉन से
साफ़ दिखाई देता बल। रोज़मर्रा का पदार्थ $10^{-9}$ से कहीं बेहतर
उदासीन है।

\textbf{11.} \omterm{def:g11:fundamental-interactions:four}{विद्युत्-चुंबकीय अंतःक्रिया}। प्रबल बल लगभग
$\qty{e-14}{m}$ से आगे \omterm{def:g11:fundamental-interactions:strong}{परास} के बाहर है --- अकेले \omterm{def:g11:fundamental-interactions:ladder}{परमाणु} से भी
$10^{4}$ गुना छोटा।

\textbf{12.} $F = \num{8.99e9} \times (\num{1.60e-19})^2 /
(\num{1.0e-15})^2 \approx \qty{2.3e2}{N}$।

\textbf{13.} $a = F/m_p = 230/\num{1.67e-27} \approx
\qty{1.4e29}{m/s^2}$, यानी $g$ का लगभग $\num{1.4e28}$ गुना: रोज़मर्रा का कोई
बल इसके पास तक नहीं फटकता।

\textbf{14.} $\qty{e-15}{m}$ पर \qty{230}{N} से अधिक प्रबल (नाभिक टिके रहते हैं);
प्रोटॉन और न्यूट्रॉन पर एक-सी लगती, आवेश से अंधी (न्यूट्रॉन भी जकड़े जाते
हैं); \omterm{def:g11:fundamental-interactions:strong}{परास} लगभग $\qty{e-15}{m}$ (नाभिक छोटे ही रहते हैं --- गोंद इससे आगे
पहुँच ही नहीं पाती)।

\textbf{15.} $F = \num{2.30e-28}/(\num{1.4e-14})^2 \approx
\qty{1.2}{N}$। नहीं: \omterm{def:g11:fundamental-interactions:strong}{परास} का $14$ गुना। प्रतिकर्षण
\emph{सभी} $\frac{92\times91}{2} = 4186$ प्रोटॉन-युग्मों पर लगता है; गोंद केवल निकटतम पड़ोसियों को
बाँधती है, इसलिए बढ़ने का फ़ायदा प्रतिकर्षण को मिलता है। न्यूट्रॉन बिना
प्रतिकर्षण जोड़े गोंद और फ़ासला दोनों देते हैं --- पर लगभग $80$ प्रोटॉन
से आगे कोई मिश्रण संतुलन नहीं बिठाता और नाभिक क्षय कर जाता है।

\textbf{16.} \omterm{def:g11:fundamental-interactions:weak}{दुर्बल अंतःक्रिया} न्यूट्रॉन को प्रोटॉन में बदल देती है
(और उलटा भी), जिससे $\beta$ क्षय होता है --- वह रूपांतरण करती है, बाँधती
कभी नहीं।

\textbf{17.} $F = \num{6.67e-11} \times \num{5.97e24} \times
\num{7.35e22}/(\num{3.84e8})^2 \approx \qty{2.0e20}{N}$।

\textbf{18.} $Q_E = \num{1.8e51} \times \num{1.6e-19} =
\qty{2.9e32}{C}$, $Q_M = \qty{3.5e30}{C}$। बलों को बराबर करने का अर्थ है
$f^2\,k\,Q_E Q_M = G M_E M_M$:
\[
f = \sqrt{\frac{\num{2.93e37}}{\num{8.99e9} \times \num{2.9e32}
\times \num{3.5e30}}} = \sqrt{\num{3.2e-36}} \approx \num{1.8e-18} :
\]
यानी \emph{$10^{18}$ में दो भाग} का असंतुलन गुरुत्व को काट देता।

\textbf{19.} आवेश दो चिह्नों में आता है और कट जाता है (प्रश्न 18 दिखाता है
कि कितनी पूर्णता से), इसलिए बड़े पैमानों पर प्रबल बल ख़ुद को बंद कर
लेता है; द्रव्यमान का एक ही चिह्न है, वह सदा जुड़ता जाता है, और उसे कोई ढाल
नहीं ढकती।

\textbf{20.} नाभिक: $\qty{e-15}{m}$ पर \omterm{def:g11:fundamental-interactions:strong}{प्रबल अंतःक्रिया} प्रोटॉनों
के प्रतिकर्षण से ज़ोर में जीत जाती है। \omterm{def:g11:fundamental-interactions:ladder}{परमाणु} से पहाड़ तक:
\omterm{def:g11:fundamental-interactions:four}{विद्युत्-चुंबकीय अंतःक्रिया} इलेक्ट्रॉनों को नाभिकों से और
परमाणुओं को एक-दूसरे से बाँधती है। ग्रह और उससे ऊपर: आवेश कट जाने
के बाद निर्विरोध \omterm{def:g11:fundamental-interactions:four}{गुरुत्वाकर्षण} ग्रहों, कक्षाओं और
आकाशगंगाओं को थामे रखता है। दुनिया भरभराकर इसलिए नहीं गिरती कि हर मंज़िल को
बोझ से सख़्त कोई बंधक टेके हुए है, और बिखरती इसलिए नहीं कि पदार्थ लगभग
$10^{-18}$ तक उदासीन है --- सो हर पैमाने पर ठीक एक बल की हुकूमत है,
और जहाँ आकर्षी होना चाहिए वहाँ वह आकर्षी है।
\end{solution}
